\begin{trapbox}
	\begin{description}[style=nextline, leftmargin=2.8em, labelsep=0.5em, itemsep=0.35em]
		\item[\textbf{\textcolor{CTrap}{易错点一（区间换元遗忘）}}] 学生直接代入 $x$ 区间到 $y=A\sin(\omega x+\varphi)$，忽略先换元，导致无法正确求出最值。
		\item[\textbf{\textcolor{CTrap}{正确理解（先换元再求值）：}}] 必须令 $t=\omega x+\varphi$，先求 $t$ 的范围，然后转化为 $\sin t$ 在 $t$ 区间上的最值问题。
		\item[\textbf{\textcolor{CTrap}{错因分析（复合处理缺失）：}}] 未识别出函数是复合函数，缺乏换元意识，胡乱代入端点。
	\end{description}

	\medskip
	\begin{description}[style=nextline, leftmargin=2.8em, labelsep=0.5em, itemsep=0.35em]
		\item[\textbf{\textcolor{CTrap}{易错点二（正弦最值误判）}}] 误认为 $\sin t$ 在 $[t_1,t_2]$ 上的最值总是 $\pm1$，不根据区间判断实际最值。
		\item[\textbf{\textcolor{CTrap}{正确理解（图象定位最值）：}}] 应在 $t$ 区间内画出 $\sin t$ 的图象或使用单位圆，确定该区间内的最大值点与最小值点。
		\item[\textbf{\textcolor{CTrap}{错因分析（忽视区间限制）：}}] 忘记正弦函数在区间上可能取不到 $\pm1$，需要具体分析。
	\end{description}

	\medskip
	\begin{description}[style=nextline, leftmargin=2.8em, labelsep=0.5em, itemsep=0.35em]
		\item[\textbf{\textcolor{CTrap}{易错点三（系数漏乘结果）}}] 求出 $\sin t$ 的最值后，直接将其作为原函数最值，忘了乘以系数 $A$（若有常数项还忘了加减常数）。
		\item[\textbf{\textcolor{CTrap}{正确理解（乘A再加常数）：}}] 最终原函数的最值应为 $A$ 乘正弦最值，再加减常数项（如果函数有常数部分）。
		\item[\textbf{\textcolor{CTrap}{错因分析（外层函数遗忘）：}}] 只计算了内层函数的值域，忽略了外层系数（和常数）对最值的缩放作用。
	\end{description}
\end{trapbox}
